\documentclass[10pt]{mypackage}

\usepackage[uprightgreeks,noDcommand]{kpfonts}
\renewcommand{\coloneq}{\coloneqq}
\renewcommand{\eqcolon}{\eqqcolon}
%\usepackage{dsfont}
%\renewcommand{\mathbb}{\mathds}
%\usepackage{homework}
%\usepackage{notes}

%\usepackage[ backend=bibtex, style = alphabetic, sorting=ynt ]{biblatex}
%\addbibresource{  }

\usepackage{parskip}

\fancyhf{}
\fancyhead[R]{Avinash Iyer}
\fancyhead[L]{Functional Analysis: Homework 1}
\fancyfoot[C]{\thepage}

\setcounter{secnumdepth}{0}

\begin{document}
\RaggedRight
\begin{problem}[Problem II.1]
  Let $V$ and $W$ be $\K$-TVSs, and let $T\colon V\rightarrow W$ be a continuous linear map. Show that if $\left( x_i \right)_i$ is a Cauchy net in $V$, then $\left( Tx_i \right)_i$ is a Cauchy net in $W$.
\end{problem}
\begin{solution}
  Let $U$ be a neighborhood of $0$ in $W$. Then, the inverse image $T^{-1}\left( U \right)$ is an open neighborhood of $0$ in $V$. There is some index $k$ such that for all $i,j\in I$ with $k\preceq i$ and $k\preceq j$, we have that $x_i-x_j\in T^{-1}(U)$. Yet, this means that $Tx_i - Tx_j\in U$ for all $i,j\in I$ with $k\preceq i$ and $k\preceq j$. This gives the desired result.
\end{solution}
\begin{problem}[Problem II.2]
  Let $\left( V, \mathcal{T}\left( \mathcal{N} \right) \right)$ be a LCTVS, and let $\left( x_i \right)_i$ be a net in $V$. Let $ \overline{x}\in V $ be some element. Show that the following are equivalent:
  \begin{enumerate}[(i)]
    \item the net $\left( x_i \right)_i$ converges to $ \overline{x} $;
    \item for all $\rho\in \mathcal{N}$ and for all $\ve > 0$, there exists $k\in I$ such that for all $i\in I$ with $k\preceq i$, one has $\rho\left( x_i - \overline{x} \right) \leq \ve$;
    \item for all continuous seminorms $\rho$ and for all $\ve > 0$, there exists $k\in I$ such that for all $i\in I$ with $k\preceq i$, one has $\rho\left( x_i - \overline{x} \right) \leq \ve$.
  \end{enumerate}
\end{problem}
\begin{solution}
  Since translation by a vector is a homeomorphism, we will assume that $ \overline{x} = 0 $ and prove the results accordingly.

  Suppose $\left( x_i \right)_i\rightarrow 0$. If $\rho$ is a defining seminorm, then we observe that the subbasic open set
  \begin{align*}
    U\left( 0,\rho,\ve \right) &= \set{y\in V : \rho(y) < \ve}
  \end{align*}
  is an open neighborhood of $0$. In particular, this means there is some $k\in I$ such that for all $i\in I$ with $k\preceq i$, one has $x_i\in U\left( 0,\rho,\ve \right)$, or that $\rho\left( x_i \right) < \ve$. This gives (ii).

  To show (iii) from (ii), one observes that if $\rho$ is a continuous seminorm, then there are defining seminorms $\rho_1,\dots,\rho_n$ and a constant $C$ such that
  \begin{align*}
    \rho\left( x \right) &\leq C \max\set{\rho_1(x),\dots,\rho_n(x)}.
  \end{align*}
  Given $\ve > 0$, for each $j=1,\dots,n$ we select $k_j\in I$ such that for all $i_j\in I$ satisfying $k_j\preceq i_j$, one has $\rho\left( x_{i_j} \right) < \ve/C$. There is some $k$ such that $k_j\preceq k$ for each $j=1,\dots,n$. Then, for any $i\in I$ with $k\preceq i$, one has $\rho\left(x_i\right)\leq \ve$, yielding (iii).
\end{solution}
\end{document}
